\lysh{14702}{4}

\begin{alist}
\item Για τη περίμετρο της μακέτας έχουμε:
$\Pi(x) = 2x + 2(2x - 1) = 6x - 2$, με $x > 2$

Και για το εμβαδόν της:
$E(x) = x(2x - 1) = 2x^2 - x$, με $x > 2$.

\item Στη περίφραξη του πάρκου εμπλέκεται η περίμετρος που δίνεται από τον τύπο
$\Pi(x) = 6x - 2$. Αρκεί $\Pi(x) \le 8$ ή $6x - 2 \le 8 \Leftrightarrow 6x \le 10 \Leftrightarrow x \le \dfrac{10}{6} \Leftrightarrow x \le \dfrac{5}{3}$.
Και για το $2x - 1$ προκύπτει $2 \cdot x \le 2 \cdot \dfrac{5}{3} \Leftrightarrow 2x \le \dfrac{10}{3} \Leftrightarrow 2x - 1 \le \dfrac{10}{3} - 1 \Leftrightarrow 2x - 1 \le \dfrac{7}{3}$.

Επειδή οι διαστάσεις παίρνουν μόνο θετικές τιμές, οι τιμές τους κυμαίνονται:
$0 < x \le \dfrac{5}{3}$ και $0 < 2x - 1 \le \dfrac{7}{3}$.

\item Το εμβαδόν της μακέτας δίνεται από το τύπο $E(x) = 2x^2 - x$.
Αρκεί $E(x) \le 1$ ή $2x^2 - x \le 1$ τότε $2x^2 - x - 1 \le 0 \quad (1)$.

Παρατηρούμε πως το τριώνυμο της ανίσωσης $(1)$ μηδενίζεται για $x_1 = 1$ και $x_2 = -\dfrac{1}{2}$ και
επειδή το $\alpha=2 >0$ προκύπτει:

\begin{center}
\begin{tabular}{|c|ccccccc|}

$x$ & $-\infty$ & & $-\dfrac{1}{2}$ & & $1$ & & $+\infty$ \\ 
$2x^2 - x - 1$ & & $+$ & $0$ & $-$ & $0$ & $+$ & \\ 
\end{tabular}
\end{center}

Σύμφωνα με τον πίνακα προσήμων το τριώνυμο θα είναι ετερόσημο του $\alpha$, δηλαδή αρνητικό
για $-\dfrac{1}{2} \le x \le 1$.
Άρα οι λύσεις της ανίσωσης $(1)$ είναι $x \in \left[-\dfrac{1}{2}, 1\right]$.
Επειδή απαιτείται $x > \dfrac{1}{2}$ ώστε να είναι $2x - 1 > 0$, τότε $x \in \left(\dfrac{1}{2}, 1\right]$.
Σχόλιο: Η απαίτηση $x>2$ στο $\alpha)$ ερώτημα δημιουργεί πρόβλημα διότι τα ερωτήματα β) και γ) είναι αδύνατα. Υποθέτουμε ότι η αρχική απαίτηση ήταν $x> \dfrac{1}{2}$.
\end{alist}
